\chapter{Обсуждение результатов}
\label{ch:discussion}

Помимо достижения целевых показателей, в работе получены три результата,
касающиеся самого измерения: определена область, в которой преимущество
планировщика исчезает; показано, что серверный показатель доступности
федерации не видит отказа кластера; обнаружено, что механизмы
самовосстановления Kubernetes могут усиливать сетевой отказ.

\section{Отрицательный результат: устойчивость к весам}

Гипотеза H5 была сформулирована строго~--- «во всех точках сетки»~--- и не
подтвердилась в одной точке из 65. В этой точке 0.9 веса приходится на
загрузку, а задержка не учитывается. Причина понятна: на стенде все узлы
работают на одной физической машине, загрузка почти не зависит от
распределения нагрузки и не несёт информации о сетевой деградации. Вдоль
всей границы $w_2 = 0$ доля размещений в деградировавший кластер растёт по
мере уменьшения веса доли ошибок, поскольку доля ошибок при сетевой
деградации нарастает медленно. Отсюда следует, что метод работает, пока
в оценке присутствует метрика, чувствительная к типу отказа: для сетевой
деградации это задержка. Заранее зафиксированные веса $(0.5, 0.3, 0.2)$
находятся далеко от проблемной области.

\section{Серверный показатель доступности не видит отказа кластера}

Первоначально глобальная доступность федерации вычислялась как сумма
серверных метрик frontend всех кластеров. При отказе кластера целиком его
запросы завершаются ошибкой на шлюзе и до frontend не доходят, то есть
выпадают из знаменателя. В прогонах без переключения трафика клиент
наблюдал 49.9\% ошибок, шлюз~--- 47.5\%, а серверный показатель всё это
время показывал 0\% ошибок.

Это прямой ответ на вопрос об измерении надёжности мультикластерного
сервиса: агрегация серверных показателей кластеров слепа к отказу
кластера, и доступность федерации необходимо измерять в точке входа
трафика. Показатель перенесён на шлюз, серверный ряд оставлен для
диагностики. Проверка на выгруженных рядах шлюза показала, что скорость
расходования бюджета ошибок за 5~минут при отказе достигает 49, то есть
бюджет ошибок расходуется в 49 раз быстрее допустимого для цели 99\%.

\section{Самовосстановление усиливает сетевой отказ}

В первой серии экспериментов сценария~2 проверки работоспособности backend
выполнялись с таймаутом 1~с. При задержке 500~мс и потере 10\% пакетов
kubelet принимал медленный, но работающий backend за неработающий и
перезапускал его. В результате backend был недоступен от 24\% до 83\%
времени отказа и восстанавливался до 159~с после снятия отказа: сетевая
деградация превращалась в отказ сервиса.

Проверки были исправлены с учётом допустимой задержки сети, и серия
повторена; приёмка гипотез выполнена по повторной серии. Первая серия
сохранена как самостоятельный результат: в ней доля ошибок у наивного
планировщика втрое выше (16.34\% против 4.41\%), а снижение у
разрабатываемого~--- больше (64.4\%). Практический вывод: таймауты проверок
работоспособности должны учитывать допустимую задержку сети, иначе
самовосстановление само становится источником отказа.

\section{Почему не сработал алерт по бюджету ошибок}

Несрабатывание алерта является свойством многооконных правил~\cite{sre-workbook},
а не дефектом. Быстрый алерт требует более 14.4\% ошибок за час, то есть при
50 запросах в секунду около 23~тысяч ошибок. Наиболее тяжёлый инцидент на
стенде (отказ без переключения трафика) дал до 7025 ошибок и 0.60\% месячного
бюджета. Для коротких инцидентов быстрее срабатывает не алерт по бюджету
ошибок, а собственный мониторинг контроллера восстановления (TTD 11~с).

\section{Ограничения и угрозы достоверности}

\begin{itemize}
  \item \textbf{Одна машина.} Абсолютные значения (задержки, TTD) зависят от
        стенда и на реальные системы не переносятся. Сравнения выполнены в
        одинаковых условиях и переносятся как направление эффекта.
  \item \textbf{Два кластера федерации.} После деградации одного из них у
        планировщика остаётся единственный кандидат, а результат получен
        для одного значения числа кластеров.
  \item \textbf{Три повтора без статистического теста.} Это компенсируется
        требованием одинакового знака эффекта в каждой паре: снижение доли
        ошибок по парам составило от 39\% до 51\%, потерь запросов~--- от
        86.8\% до 87.3\%.
  \item \textbf{Малоинформативная загрузка.} Узлы делят одну машину, и вклад
        третьего слагаемого оценки мал.
  \item \textbf{Нейтральная оценка кластера без трафика.} Кластер без
        метрик получает оценку около 0.55 против около 0.95 у кластера с
        трафиком, поэтому с холодного старта планировщик собирает поток в
        одном кластере. На стенде это обойдено фоновой нагрузкой; в формуле
        это ограничение сохраняется.
  \item \textbf{Только новые размещения.} Планировщик не переносит уже
        размещённые нагрузки, поэтому его эффект ограничен скоростью
        обновления потока.
  \item \textbf{Единая точка отказа прототипа.} Управляющий кластер и шлюз
        не дублируются; в реальной системе их резервирование обязательно,
        но в задачу работы оно не входило.
  \item \textbf{Нагрузка без состояния.} Репликация данных и допустимая
        потеря данных не рассматривались.
\end{itemize}
